\chapter{Executive Summary}
\label{ch:executive-summary}

\section{What the product is}

Expo Snips is a free developer tool that writes complete Expo / React Native
screens into an existing mobile project from a single short command or an editor
trigger. A developer types \texttt{expo-snips add auth/login} in a terminal, or
\texttt{snip:} inside VS Code, and the tool produces the screen file, its
supporting form component and any hooks --- complete with form validation,
loading and error states, and keyboard handling.

\textbf{Status:} IMPLEMENTED.

\section{The problem}

Teams building React Native apps repeatedly hand-write the same screens:
login, registration, one-time-code entry, settings, product listings, onboarding
flows and empty/error states. Each one involves several files and several
libraries, and it is easy to omit the details that make a screen feel finished
(validation messages, loading spinners, keyboard avoidance, safe areas).

\textbf{Status:} INFERRED from the product's stated rationale
(\texttt{packages/vscode-extension/README.md}) and from the recipe coverage in
\texttt{packages/registry/data}.

\section{The solution}

A catalog of 24 versioned ``screen recipes''. Each recipe declares the stacks it
supports, the npm dependencies it needs, and the files it produces. A shared
Handlebars engine renders those files to match the target project.

\textbf{Status:} IMPLEMENTED.

\section{How value is delivered}

\begin{enumerate}
  \item \textbf{Whole screens, not snippets.} A recipe such as \texttt{auth/login}
        emits up to seven file variants depending on the project's stack
        (\texttt{app/login/index.tsx}, \texttt{screens/LoginScreen.tsx},
        \texttt{components/LoginForm.tsx}, \texttt{hooks/useLogin.ts} or
        \texttt{.js}).
  \item \textbf{Stack-aware output.} The tool reads \texttt{package.json} and
        \texttt{tsconfig.json} to detect the router, styling solution, UI
        primitives, icon library and language, then renders conditionally.
  \item \textbf{Best-practice patterns.} Auth recipes use
        \texttt{react-hook-form} with \texttt{zod} validation and
        \texttt{@hookform/resolvers}.
  \item \textbf{No dependency drift.} Recipe dependencies are installed with the
        project's own package manager (npm, pnpm, yarn or bun).
  \item \textbf{Reviewable updates.} Installs are recorded in
        \texttt{.expo-snips.json}; \texttt{expo-snips update} shows a line-level
        diff before applying changes.
  \item \textbf{Works offline.} The catalog is bundled into both the CLI and the
        extension, so no network access is required.
\end{enumerate}

\textbf{Status:} IMPLEMENTED --- each item is traceable to source (see the
corresponding chapters).

\section{Technology and architecture in one paragraph}

Expo Snips is a pnpm workspace monorepo with four packages: a shared generation
engine (\texttt{@expo-snips/core}, Handlebars + TypeScript types), a CLI
(\texttt{@expo-snips/cli}, commander + prompts + chalk), a registry package
(\texttt{@expo-snips/registry}, recipe data + 48 templates + a JSON Schema) and
a VS Code extension (\texttt{expo-snippets}, esbuild-bundled). The catalog is
compiled from source recipe files into a static JSON registry that is published
to GitHub Pages, and into offline bundles embedded in the CLI and the extension.
There is no backend, no database and no user account system.

\textbf{Status:} IMPLEMENTED.

\section{Current maturity}

\textbf{Classification: MVP / early-stage.}

\begin{itemize}
  \item All packages are version 0.1.0 (extension 0.1.3).
  \item The generation engine, CLI, registry and VS Code extension are all
        implemented and functional.
  \item Infrastructure for distribution exists (GitHub Actions for registry
        deployment and npm publishing).
  \item There are no automated tests, no linters and no root README.
  \item There is no evidence of production adoption, usage metrics, customers or
        testimonials anywhere in the repository.
\end{itemize}

\textbf{Status:} IMPLEMENTED (versions, CI, absence of tests);
\textbf{UNKNOWN} (adoption --- not evidenced).

\section{Key limitations}

\begin{itemize}
  \item Generated screens target \textbf{Expo Router} by default; only
        \texttt{auth/login} and \texttt{auth/register} support
        react-navigation, and only in their native-UI form.
  \item \texttt{tamagui}, \texttt{vector-icons} and \texttt{custom} UI are
        declared in types, schema and VS Code settings but no recipe implements
        them.
  \item Generated authentication hooks contain \emph{placeholder} logic (a
        1.5-second mock delay and an alert) and must be replaced by the developer.
  \item The shadcn-style templates rely on design tokens
        (\texttt{bg-background}, \texttt{text-foreground}, \dots) for which no
        Tailwind / NativeWind theme configuration is shipped in the repository.
  \item The VS Code extension exposes no update or status command; those exist
        only in the CLI.
\end{itemize}

\textbf{Status:} IMPLEMENTED (each limitation is cited in
Chapter~\ref{ch:known-issues}).

\section{Summary judgement}

Expo Snips is a coherent, well-factored developer tool with a genuinely useful
core idea --- stack-aware whole-screen generation from a versioned registry ---
implemented across two delivery surfaces. Its principal risks are not
architectural but of completeness and hardening: no tests, several declared but
unimplemented options, and placeholder business logic in generated code that must
be clearly communicated to users.

\textbf{Status:} INFERRED (an assessment by this document, based on the evidence
gathered).
